% Generated draft: requires mathematical and source review.
\documentclass[11pt]{article}
\usepackage[margin=25mm]{geometry}
\usepackage{fontspec}
\setmainfont{Latin Modern Roman}
\usepackage{amsmath,amssymb,mathtools,unicode-math}
\setmathfont{Latin Modern Math}
\usepackage{enumitem,needspace,xurl,hyperref}
\setlength{\parindent}{0pt}
\setlength{\parskip}{5pt}
\setlength{\emergencystretch}{3em}
\newcommand{\N}{\mathbb{N}}
\newcommand{\Z}{\mathbb{Z}}
\newcommand{\Q}{\mathbb{Q}}
\newcommand{\R}{\mathbb{R}}
\newcommand{\C}{\mathbb{C}}
\newcommand{\F}{\mathbb{F}}
\newcommand{\A}{\mathbb{A}}
\newcommand{\PP}{\mathbb{P}}
\newcommand{\E}{\mathbb{E}}
\newcommand{\eps}{\varepsilon}
\newcommand{\dd}{\,\mathrm{d}}
\begin{document}
\section*{UnsolvedMath ID: 30005717}
\textbf{Source problem code:} catalogue entry
\subsection{Definitions and mathematical statement}
Let \(k\in\Z\) with \(k\geq2\), and put \(d=2k\). A simplicial \(d\)-polytope is a full-dimensional convex polytope \(P\subset\R^d\) whose proper faces are simplices. Its natural embedding places each vertex at its actual position in \(\R^d\). Write \(V(P)\) for its vertex set and \(\Delta(P)\) for its boundary complex. A missing face is a nonempty subset \(F\subseteq V(P)\) such that \(F\notin\Delta(P)\) but every proper subset of \(F\) belongs to \(\Delta(P)\); its dimension is \(|F|-1\). The supplied hypothesis is
\[
F\text{ a missing face of }\Delta(P)\quad\Longrightarrow\quad |F|-1<d-k+1=k+1.
\]
Thus every missing face has dimension at most \(k\).

Write \(\mathcal S_k^{\mathrm{aff}}(P)\) for the space termed the \emph{space of affine \(k\)-stresses} in the supplied record. The intended assertion is that, for every such \(P\), this stress-space data determines \(P\) up to affine equivalence. Here \(P\) and \(Q\) are affinely equivalent if there exist \(A\in\operatorname{GL}_d(\R)\) and \(b\in\R^d\) such that \(Q=\{Ax+b:x\in P\}\).

For a comparison formulation within the stated class, let \(P,Q\subset\R^{2k}\) satisfy the preceding hypotheses and let \(\phi:V(P)\to V(Q)\) be a bijection. The intended reconstruction implication is that identification of their affine \(k\)-stress data through \(\phi\) forces affine equivalence of \(P\) and \(Q\). However, the supplied record does not define affine \(k\)-stresses, their ambient coordinate space, or the data retained when two stress spaces are identified. It also does not specify whether comparison polytopes must satisfy the missing-face hypothesis or whether the resulting affine equivalence must realize \(\phi\). Consequently, this comparison formulation is conditional on resolving those conventions; an abstract vector-space isomorphism is not silently adopted as the meaning of identical stress data.
\subsection{Short English statement}
For each integer \(k\geq2\), the supplied conjectural assertion says that a naturally embedded simplicial \(2k\)-polytope with no missing faces of dimension at least \(k+1\) can be reconstructed up to affine equivalence from its affine \(k\)-stress space. The precise stress-space data and comparison conventions require clarification from the cited source.
\subsection{Sources}
\begin{itemize}
\item Supplied local registry record, permanent UnsolvedMath ID \(30005717\), title \emph{Even-Dimensional Stress-Space Reconstruction}, exact source problem code \texttt{catalogue entry}. The registry records proposed year \(2023\).
\item \emph{Geometric, Algebraic and Topological Combinatorics} (2024), Oberwolfach Reports, pp.\ 3249--3318, DOI \texttt{10.4171/owr/2023/58}: \url{https://doi.org/10.4171/owr/2023/58}. The supplied original excerpt carries page number \(3303\), attributes the surrounding conjecture to Kalai through reference [4], and states that the \(d=2k\geq4\) case remains open. The bibliographic identity of reference [4] is not supplied.
\item Supplied literature source labelled ``Novik--Zheng, affine stresses and reconstruction'': \url{https://arxiv.org/abs/2208.06693}.
\item Literature assessment dated exactly \texttt{2026-08-22}: the supplied curation labels the broader problem ``partially solved'', reports that the \(k=1\) and \(k=\lfloor d/2\rfloor+1\) boundary cases are settled, and describes the even-dimensional \(d=2k\geq4\) case as still open. This is supplied dated curation, not an independently verified literature assessment; the boundary cases concern the broader conjecture rather than additional cases within the restricted parameters above.
\item The registry labels the statement \texttt{corrected\_verified} and says it was reconstructed and checked against the official report. These are supplied provenance claims, not verification performed for this draft.
\end{itemize}
\end{document}
